%%  JAAMAS asks for 150 to 250 words, no citations, no equations, no undefined
%%  abbreviations, and 4 to 6 keywords.  build.py counts both.
\abstract{An ecosystem of deployed autonomous agents is selected twice over. An
agent earns a payoff by acting for a principal, but whether its design spreads
is decided by that principal, who never plays the game and bears only part of
the harm it causes. Evolutionary game theory normally identifies the two. We
call their difference the selection wedge, embed it in a four-strategy model of
an artificial intelligence race, and analyse its replicator and
finite-population dynamics. The wedge collapses to a single control parameter,
an effective liability, in which every invasion condition is affine, so all
thresholds are closed form. Four consequences follow for governance. Screening
designs before deployment, at any tolerance that admits reciprocation, removes
only a weakly dominated design and leaves the long-run unsafe frequency
unchanged. That frequency is not monotone in liability: over an intermediate
window liability also taxes the retaliation on which conditional safety
depends, the unsafe state becomes uninvadable by conditional safety, and
crossing the window's lower edge from below lifts the basin-weighted unsafe
frequency from zero to about a third. The barrier protecting a safe population
meanwhile varies by a factor of 77.7, at the baseline calibration, along a
composition that selection leaves neutral. And irreversible capability
diffusion makes the loss of a protected regime hysteretic, with exactly
computable width. Governance of such an ecosystem binds at the layer that
selects designs, not at the layer that plays them.}

\keywords{Multi-agent systems, Evolutionary game theory, Replicator dynamics,
Liability, Artificial intelligence governance}
